\documentclass[12pt,a4paper,openany,oneside]{book}

% --- Paquetes ---
\usepackage[utf8]{inputenc}
\usepackage[spanish, es-noshorthands]{babel}
\usepackage{amsmath, amsfonts, amssymb}
\usepackage{graphicx}
\usepackage{geometry}
\usepackage{hyperref}
\usepackage{booktabs}
\usepackage{fancyhdr}
\usepackage{listings}
\usepackage{xcolor}
\usepackage{tikz}
\usepackage{pgfplots}
\usepackage{colortbl}
\usepackage{multirow}
\usepackage{hhline}
\usetikzlibrary{arrows.meta, positioning, shapes.geometric, calc}
\pgfplotsset{compat=1.18}

\geometry{margin=2.5cm}

% --- Colores y estilos para código Python ---
\definecolor{codegreen}{rgb}{0,0.6,0}
\definecolor{codegray}{rgb}{0.5,0.5,0.5}
\definecolor{codepurple}{rgb}{0.58,0,0.82}
\definecolor{backcolour}{rgb}{0.95,0.95,0.92}

\lstset{
    language=Python,
    backgroundcolor=\color{backcolour},   
    commentstyle=\color{codegreen},
    keywordstyle=\color{blue},
    numberstyle=\tiny\color{codegray},
    stringstyle=\color{codepurple},
    basicstyle=\ttfamily\small,
    breaklines=true,
    frame=single,
    showstringspaces=false
}

% --- Estilo de cabecera y pie ---
\pagestyle{fancy}
\fancyhf{}
\renewcommand{\headrulewidth}{0.4pt}
\renewcommand{\footrulewidth}{0.4pt}
\fancyhead[L]{\small Carlos César Sánchez Coronel}
\fancyhead[R]{\small Sesión 10: Series de Tiempo}
\fancyfoot[L]{\small Carlos César Sánchez Coronel}
\fancyfoot[C]{\thepage}

% --- Portada ---
\title{
    \textbf{Sesión 10} \\ 
    \Huge \textbf{SERIES DE TIEMPO} \\
    \Large Modelado estadístico y aprendizaje automático para datos temporales
}
\author{
    \small Por \\ \vspace{0.2cm}
    \Large \textsc{Carlos César Sánchez Coronel}
}
\date{2026}

\begin{document}

\maketitle
\tableofcontents
\addcontentsline{toc}{chapter}{Índice general}

% ------------------------------------------------------------------
% CAPÍTULO 10: SESIÓN 10 - SERIES DE TIEMPO
% ------------------------------------------------------------------
\chapter{Series de Tiempo}

Los datos temporales aparecen en innumerables contextos: ventas diarias, tráfico web, indicadores económicos, sensores IoT, etc. Modelar correctamente la dependencia temporal es esencial para realizar pronósticos precisos y tomar decisiones basadas en datos. Esta sesión aborda tanto los modelos estadísticos clásicos (ARIMA, SARIMA) como los enfoques modernos de machine learning (XGBoost, Prophet), junto con las técnicas de feature engineering específicas para series temporales. Se presentan los fundamentos matemáticos, las métricas de evaluación y ejemplos prácticos en cada caso.

\section{Logro de la sesión}
Modelar datos temporales utilizando técnicas estadísticas y de machine learning, comprendiendo las particularidades de este tipo de datos y las métricas adecuadas para evaluar pronósticos.

\section{Componentes de una serie temporal}

Una serie temporal \(y_t\) (con \(t=1,\dots,T\)) puede descomponerse en varios componentes no observables directamente:

\begin{itemize}
    \item \textbf{Tendencia (Trend)}: movimiento a largo plazo, ya sea creciente, decreciente o estable.
    \item \textbf{Estacionalidad (Seasonal)}: patrones regulares que se repiten en períodos fijos (diario, semanal, anual).
    \item \textbf{Ciclo (Cycle)}: fluctuaciones sin período fijo, asociadas a condiciones económicas o de negocio (ej. ciclos de auge y recesión). A diferencia de la estacionalidad, el ciclo no tiene una duración predecible.
    \item \textbf{Ruido (Noise)}: variaciones aleatorias no explicadas por los demás componentes.
\end{itemize}

\subsection{Descomposición clásica}
Existen dos formas principales de combinar estos componentes:

\begin{itemize}
    \item \textbf{Aditiva}: \(y_t = T_t + S_t + C_t + R_t\)
    \item \textbf{Multiplicativa}: \(y_t = T_t \times S_t \times C_t \times R_t\)
\end{itemize}

La descomposición multiplicativa es apropiada cuando la magnitud de las fluctuaciones estacionales o aleatorias crece con el nivel de la serie (por ejemplo, ventas que aumentan en Navidad cada año de forma proporcional al nivel).

\subsection{Descomposición en la práctica}
Métodos como STL (Seasonal and Trend decomposition using Loess) permiten estimar estos componentes de forma robusta. La descomposición es útil para:
\begin{itemize}
    \item Entender el comportamiento de la serie.
    \item Eliminar la estacionalidad (ajuste estacional).
    \item Crear características para modelos de machine learning.
\end{itemize}

\section{Feature engineering para series temporales}

Cuando se utilizan modelos de machine learning estándar (regresión, árboles, XGBoost) con datos temporales, es necesario convertir el problema de series de tiempo en un problema de aprendizaje supervisado mediante la creación de características a partir de la historia.

\subsection{Rezagos (lags)}
La característica más básica es el valor de la serie en instantes anteriores:
\[
y_{t-1}, y_{t-2}, \dots, y_{t-L}
\]
donde \(L\) es el número máximo de rezagos. El modelo aprende la relación entre el valor actual y los pasados.

\subsection{Ventanas móviles}
Estadísticas calculadas sobre una ventana de tiempo:
\begin{itemize}
    \item Media móvil: \(\frac{1}{w}\sum_{i=t-w}^{t-1} y_i\)
    \item Desviación estándar móvil
    \item Mínimo, máximo, percentiles
\end{itemize}
Estas características capturan la tendencia local y la volatilidad.

\subsection{Diferencias}
Para hacer la serie estacionaria (ver sección ARIMA), se suele trabajar con diferencias:
\[
\Delta y_t = y_t - y_{t-1}
\]
Las diferencias también pueden incluirse como características.

\subsection{Variables temporales (calendario)}
\begin{itemize}
    \item Hora del día, día de la semana, mes, año.
    \item Indicadores de festivos, fines de semana, eventos especiales.
    \item Variables cíclicas: \(\sin(2\pi \cdot \text{hora}/24)\), \(\cos(2\pi \cdot \text{hora}/24)\) para preservar la naturaleza circular del tiempo.
\end{itemize}

\subsection{Agregaciones en múltiples escalas}
Por ejemplo, para predicción diaria, se pueden incluir:
\begin{itemize}
    \item Media de los últimos 7 días.
    \item Media del mismo día de la semana en las últimas 4 semanas.
    \item Mismo día del año anterior (para estacionalidad anual).
\end{itemize}

\section{Modelos clásicos: ARIMA y SARIMA}

Los modelos ARIMA (AutoRegressive Integrated Moving Average) son la piedra angular del análisis de series temporales.

\subsection{ARMA (p,q)}
Un proceso ARMA(p,q) combina un componente autorregresivo (AR) de orden \(p\) y uno de medias móviles (MA) de orden \(q\):
\[
y_t = c + \phi_1 y_{t-1} + \dots + \phi_p y_{t-p} + \varepsilon_t + \theta_1 \varepsilon_{t-1} + \dots + \theta_q \varepsilon_{t-q}
\]
donde \(\varepsilon_t\) es ruido blanco (media cero, varianza constante).

\subsection{ARIMA(p,d,q)}
Para series no estacionarias, se aplican diferencias. El parámetro \(d\) indica el número de diferencias necesarias para lograr estacionariedad. Un ARIMA(p,d,q) equivale a un ARMA(p,q) aplicado a la serie diferenciada \(d\) veces.

\subsection{Estacionalidad: SARIMA(p,d,q)(P,D,Q)s}
Cuando hay estacionalidad con período \(s\), se incorporan componentes autorregresivos y de medias móviles estacionales. Por ejemplo, para datos mensuales con estacionalidad anual (\(s=12\)), se tiene:
\[
\Phi(B^s) \phi(B) (1-B^s)^D (1-B)^d y_t = \Theta(B^s) \theta(B) \varepsilon_t
\]
donde \(B\) es el operador de rezago (\(B y_t = y_{t-1}\)), \(\phi\) y \(\theta\) son polinomios regulares, \(\Phi\) y \(\Theta\) son polinomios estacionales.

\subsection{Selección de órdenes}
\begin{itemize}
    \item \textbf{ACF (autocorrelación)} y \textbf{PACF (autocorrelación parcial)}: ayudan a identificar \(p\) y \(q\).
    \item \textbf{Criterios de información}: AIC, BIC para comparar modelos.
    \item \textbf{Validación}: se evalúa el modelo en una muestra de prueba (últimos períodos).
\end{itemize}

\subsection{Ejemplo práctico: predicción de pasajeros de aerolíneas}
Serie clásica con tendencia creciente y estacionalidad anual. Se aplica logaritmo para estabilizar la varianza, luego se diferencian para eliminar tendencia y estacionalidad. Un SARIMA(0,1,1)(0,1,1)\(_{12}\) suele ajustarse bien.

\section{Modelos de machine learning para series}

En los últimos años, los modelos de machine learning han ganado popularidad en pronósticos, especialmente cuando hay muchos predictores externos.

\subsection{Regresión con características temporales}
Se construye un dataset con las características descritas en la sección de feature engineering y se entrena un modelo de regresión (lineal, árboles, etc.). Este enfoque es flexible y permite incluir variables exógenas.

\subsection{XGBoost para series temporales}
XGBoost y otros gradient boosting son muy efectivos. Requieren:
\begin{itemize}
    \item Creación cuidadosa de lags y ventanas.
    \item Manejo de la estacionalidad mediante variables indicadoras.
    \item Validación cruzada temporal (no aleatoria).
\end{itemize}

\textbf{Ventajas}: capturan relaciones no lineales, manejan muchas características, robustos a outliers.

\textbf{Desventajas}: pueden tener problemas para extrapolar tendencias fuera del rango de entrenamiento.

\subsection{Prophet (Facebook)}
Prophet es un modelo aditivo desarrollado por Facebook que descompone la serie en:
\[
y(t) = g(t) + s(t) + h(t) + \varepsilon_t
\]
donde:
\begin{itemize}
    \item \(g(t)\): tendencia (lineal por tramos o logística con saturación).
    \item \(s(t)\): estacionalidad (mediante series de Fourier).
    \item \(h(t)\): efectos de días festivos (definidos por el usuario).
\end{itemize}
Prophet es robusto a datos faltantes, maneja cambios de tendencia y es fácil de usar.

\textbf{Ventajas}: automático en muchos aspectos, interpretable, maneja festivos.

\textbf{Desventajas}: menos flexible que XGBoost para incorporar muchos predictores externos.

\subsection{Comparación de enfoques}
\begin{table}[h]
\centering
\begin{tabular}{|l|c|c|c|}
\hline
\textbf{Característica} & \textbf{ARIMA/SARIMA} & \textbf{XGBoost} & \textbf{Prophet} \\
\hline
Fundamento & Estadístico & Árboles boosting & Aditivo generalizado \\
Estacionalidad & Explícita (SARIMA) & Variables dummy & Series de Fourier \\
Tendencia & Diferencias & Características & Lineal por tramos \\
Variables exógenas & Limitado (ARIMAX) & Sí & No nativo (solo festivos) \\
Extrapolación de tendencia & Sí (con tendencia) & Limitada & Sí \\
Facilidad de uso & Media & Alta & Muy alta \\
Interpretabilidad & Media (coeficientes) & Baja (importancia) & Alta \\
\hline
\end{tabular}
\caption{Comparación de enfoques para series temporales.}
\end{table}

\section{Métricas de evaluación para pronósticos}

Dado que se predicen valores continuos, las métricas de regresión son aplicables, pero con algunas particularidades.

\subsection{MAE (Mean Absolute Error)}
\[
\text{MAE} = \frac{1}{n} \sum_{t=1}^n |y_t - \hat{y}_t|
\]
Fácil de interpretar, en las mismas unidades de la serie.

\subsection{RMSE (Root Mean Squared Error)}
\[
\text{RMSE} = \sqrt{\frac{1}{n} \sum_{t=1}^n (y_t - \hat{y}_t)^2}
\]
Penaliza más los errores grandes, útil cuando los errores grandes son especialmente indeseables.

\subsection{MAPE (Mean Absolute Percentage Error)}
\[
\text{MAPE} = \frac{100}{n} \sum_{t=1}^n \left| \frac{y_t - \hat{y}_t}{y_t} \right|
\]
Se interpreta como porcentaje. Tiene problemas cuando \(y_t\) es cero o muy pequeño, y asimétrico (penaliza más la sobreestimación si \(y_t\) es pequeño).

\subsection{MASE (Mean Absolute Scaled Error)}
Propuesta por Hyndman y Koehler (2006) para superar las limitaciones del MAPE. Se define como:
\[
\text{MASE} = \frac{\frac{1}{n} \sum_{t=1}^n |y_t - \hat{y}_t|}{\frac{1}{n-1} \sum_{t=2}^n |y_t - y_{t-1}|}
\]
El denominador es el error del método ingenuo (pronóstico naive: \(\hat{y}_t = y_{t-1}\)). Un MASE \(<1\) indica que el modelo supera al naive.

\subsection{Selección de métrica según el caso}
\begin{itemize}
    \item \textbf{MAE}: si todos los errores tienen igual importancia.
    \item \textbf{RMSE}: si los errores grandes son críticos (ej. picos de demanda).
    \item \textbf{MAPE}: si se desea una medida relativa (y los valores no son cero).
    \item \textbf{MASE}: cuando se quiere comparar con un benchmark simple y es robusto a la escala.
\end{itemize}

\section{Caso integrado: Pronóstico de ventas diarias en retail}

\subsection{Problema}
Una cadena de supermercados desea predecir las ventas diarias de una tienda para las próximas 4 semanas. Se dispone de:
\begin{itemize}
    \item Historial de ventas diarias de los últimos 3 años.
    \item Calendario de días festivos y promociones.
    \item Datos meteorológicos (temperatura, lluvia).
    \item Indicadores económicos locales (opcional).
\end{itemize}

\subsection{Análisis exploratorio}
Se grafica la serie completa. Se observa:
\begin{itemize}
    \item Tendencia creciente (apertura de nuevas tiendas, inflación).
    \item Fuerte estacionalidad semanal (mayor venta los fines de semana).
    \item Picos en Navidad y otras festividades.
    \item Efecto de promociones (descuentos) que aumentan ventas.
\end{itemize}

\subsection{Preprocesamiento}
\begin{itemize}
    \item Manejo de valores faltantes (si los hay, imputación con media móvil).
    \item Creación de variables temporales: día de semana, mes, año, indicador de fin de semana, días hasta Navidad, etc.
    \item Variables de festivos: one-hot encoding para los principales festivos.
    \item Variables meteorológicas: temperatura media, precipitación (con posibles rezagos).
    \item Rezagos de ventas: ventas de los últimos 7, 14, 28 días.
    \item Media móvil de los últimos 7 días, y media del mismo día de la semana en las últimas 4 semanas.
\end{itemize}

\subsection{Modelado}

\subsubsection{ARIMA/SARIMA}
Se ajusta un SARIMA con estacionalidad semanal (\(s=7\)). Se prueba con transformación logarítmica para estabilizar la varianza. Se eligen órdenes mediante ACF/PACF y criterios de información. El modelo captura bien la estacionalidad pero no puede incorporar las promociones ni el clima directamente (aunque se podría usar ARIMAX).

\subsubsection{XGBoost}
Se entrena un XGBoost con todas las características creadas. Se utiliza validación cruzada temporal (time series split) para evitar usar datos futuros. Se ajustan hiperparámetros con Optuna. El modelo obtiene buen rendimiento, especialmente al capturar picos por promociones.

\subsubsection{Prophet}
Se ajusta Prophet con estacionalidad semanal y anual, y se añaden los días festivos como eventos. Prophet maneja automáticamente la tendencia y estacionalidad, pero no puede usar las variables meteorológicas fácilmente (se podrían añadir como regresores externos en versiones recientes).

\subsection{Evaluación}
Se comparan los modelos en el conjunto de prueba (últimos 28 días) usando MAE, RMSE, MAPE y MASE.

\begin{table}[h]
\centering
\begin{tabular}{l|c|c|c|c}
\textbf{Modelo} & \textbf{MAE} & \textbf{RMSE} & \textbf{MAPE (\%)} & \textbf{MASE} \\
\hline
Naive (día anterior) & 2450 & 3200 & 12.5 & 1.00 \\
SARIMA & 1850 & 2500 & 9.2 & 0.75 \\
XGBoost & 1420 & 1950 & 7.1 & 0.58 \\
Prophet & 1600 & 2100 & 8.0 & 0.65 \\
\end{tabular}
\caption{Resultados de pronóstico.}
\end{table}

XGBoost supera a los demás, seguido de Prophet. La ventaja de XGBoost se debe a su capacidad para incorporar las promociones y el clima.

\subsection{Conclusiones del caso}
El modelo XGBoost se implementa en producción, actualizándose diariamente con los nuevos datos. Se monitorea el error y se reentrena periódicamente. La empresa utiliza los pronósticos para optimizar inventarios y personal.

\section{Implementación en Python}

\subsection{ARIMA con statsmodels}
\begin{lstlisting}[language=Python]
import statsmodels.api as sm
from statsmodels.tsa.arima.model import ARIMA

model = ARIMA(series, order=(1,1,1), seasonal_order=(1,1,1,7))
fitted = model.fit()
forecast = fitted.forecast(steps=28)
\end{lstlisting}

\subsection{XGBoost con validación temporal}
\begin{lstlisting}[language=Python]
import xgboost as xgb
from sklearn.model_selection import TimeSeriesSplit

tscv = TimeSeriesSplit(n_splits=5)
for train_index, val_index in tscv.split(X):
    X_train, X_val = X[train_index], X[val_index]
    y_train, y_val = y[train_index], y[val_index]
    model = xgb.XGBRegressor(n_estimators=100, learning_rate=0.1)
    model.fit(X_train, y_train)
    # evaluar...
\end{lstlisting}

\subsection{Prophet}
\begin{lstlisting}[language=Python]
from prophet import Prophet
import pandas as pd

df = pd.DataFrame({'ds': fechas, 'y': ventas})
model = Prophet(yearly_seasonality=True, weekly_seasonality=True)
model.add_country_holidays(country_name='PE')
model.fit(df)
future = model.make_future_dataframe(periods=28)
forecast = model.predict(future)
\end{lstlisting}

\section{Conexión con el resto del curso}
\begin{itemize}
    \item El feature engineering para series utiliza conceptos de transformaciones y agregaciones vistos en sesiones anteriores.
    \item XGBoost ya se estudió en la sesión de Gradient Boosting; aquí se aplica a un dominio específico.
    \item La validación cruzada temporal complementa los métodos de validación vistos en la sesión 8.
    \item Las métricas de regresión (MAE, RMSE) ya son conocidas.
\end{itemize}

\section{Resumen}
\begin{itemize}
    \item Las series temporales se descomponen en tendencia, estacionalidad, ciclo y ruido.
    \item El feature engineering específico (rezagos, ventanas, variables de calendario) es crucial para aplicar ML.
    \item ARIMA/SARIMA son modelos estadísticos que capturan autocorrelaciones y estacionalidad de forma paramétrica.
    \item XGBoost y Prophet ofrecen alternativas flexibles, especialmente cuando hay predictores externos.
    \item La evaluación debe hacerse con métricas adecuadas (MAE, RMSE, MAPE, MASE) y mediante validación temporal.
    \item La elección del modelo depende de las características de la serie, la disponibilidad de datos externos y la necesidad de interpretabilidad.
\end{itemize}

\end{document}